\originalsection{官方 Bonus: 2011年11月29日交}
\begin{exercise}\label{ps:bonus}
\textbf{原题 I(a--i), 原卷两个f依次记为 $f_1,f_2$.} 对作业8的 $K,T,J_K$, 令
\[
\widetilde J^\mu=J_K^\mu-2t\phi\partial^\mu\phi+\phi^2\partial^\mu t.
\]
(a,b)证 $K$ 未来类时及共形 Killing 式; (c)求应力迹; (d)求 $J_K$ 散度; (e)证 $\widetilde J$ 无散度; ($f_1$)求 $J_K^0$ 零标架式; ($f_2$)求 $\widetilde J^0$; (g)证紧支撑 $\phi$ 的径向分部积分式; (h)给正平方表达; (i)证修正共形能量守恒.
\end{exercise}
\begin{solution}
(a,b,$f_1$)与作业8-II(a,b,e)完全重复, 证明见习题\ref{ps:8II}; 保留这里的课程位置而复用相同推导.
(c)4维迹为 $-Q$.
(d)对称性给 $\partial_\mu J_K^\mu=2tQ$.
(e) 在 $\Box\phi=0$ 上,
\[
\partial_\mu(-2t\phi\partial^\mu\phi)=-2\phi\partial^0\phi-2tQ,
\qquad \partial_\mu(\phi^2\partial^\mu t)=2\phi\partial^0\phi.
\]
$\partial^0=-\partial_t$, 交叉项相消, 再加(d)得0.
($f_2$) $\partial^0t=-1$, 所以 $\widetilde J^0=J_K^0+2t\phi\phi_t-\phi^2$.
(g) $\diver_x(x\phi^2)=3\phi^2+2r\phi\phi_r$. 紧支撑使散度积分0, 故 $-\int\phi^2=(2/3)\int r\phi\phi_r$. 也可用球坐标, $\int_0^\infty r^3\partial_r(\phi^2)\dd r=-3\int_0^\infty r^2\phi^2\dd r$.
(h) 置
\[
P=\tfrac14\{(L\phi)^2+[L((t+r)\phi)]^2+(\overline L\phi)^2
+[\overline L((t-r)\phi)]^2+2(1+t^2+r^2)H\}.
\]
这里 $L(t+r)=2$, $\overline L(t-r)=2$. 展开两个平方得 $P=J_K^0+2t\phi\phi_t+2r\phi\phi_r+2\phi^2$, 因而
\[
P-\widetilde J^0=2r\phi\phi_r+3\phi^2=\diver_x(x\phi^2),\qquad
\int P\dd x=\int\widetilde J^0\dd x\ge0.
\]
原卷(h)写为点态相等, 实际只有积分后相等, 或相差上述空间散度. 不把分部积分后的等式误称逐点等式.
(i) 若光滑初值紧支撑, 波方程有限传播保证每个有限时间段的统一紧支撑. 时空柱的侧面通量0, 无散度给 $\int\widetilde J^0(t)=\int\widetilde J^0(0)$. 用(h)即
\[
\mathcal E_{\rm conf}(t)=\int_{\R^3}P(t,x)\dd x=\int_{\R^3}P(0,x)\dd x.
\]
右端以 $f=\phi(0)$、$g=\phi_t(0)$ 代入 $L\phi=g+f_r$, $\overline L\phi=g-f_r$, $H=|\nabla f|^2-f_r^2$ 即由初值完全确定. 一般非紧支撑场须加有限加权能量及无穷远通量条件. 本节建立守恒, 不据此声称已证明所有零导数的精细衰减率.
\end{solution}
